\section{Conclusiones}

\begin{frame}{Conclusiones}
    \begin{enumerate}
        \setlength{\itemsep}{10pt}
        \item Un qubit basta para un motor térmico; su estadística es exacta.
        \item A máxima potencia, $\eta^* = 0{,}296 \approx \eta_{CA}$.
        \item El frente óptimo vive sobre la cota TUR ($\xi \approx 1$).
        \item MOEA/D-AWA encuentra el mejor frente; la suma ponderada no.
    \end{enumerate}
\end{frame}

\begin{frame}{Referencias}
    \footnotesize
    \begin{itemize}
        \setlength{\itemsep}{4pt}
        \item F.\,L. Curzon y B. Ahlborn, \emph{Am. J. Phys.} \textbf{43}, 22 (1975).
        \item K. Deb et al., NSGA-II, \emph{IEEE Trans. Evol. Comput.} \textbf{6}(2), 182--197 (2002).
        \item Q. Zhang y H. Li, MOEA/D, \emph{IEEE Trans. Evol. Comput.} \textbf{11}(6), 712--731 (2007).
        \item Y. Qi et al., MOEA/D-AWA, \emph{Evolutionary Computation} \textbf{22}(2), 231--264 (2014).
        \item P. Pietzonka y U. Seifert, \emph{Phys. Rev. Lett.} \textbf{120}, 190602 (2018).
        \item P.\,A. Erdman y F. Noé, \emph{npj Quantum Inf.} \textbf{8}, 1 (2022).
        \item P.\,A. Erdman et al., \emph{Phys. Rev. Research} \textbf{5}, L022017 (2023).
    \end{itemize}
\end{frame}

{
\setbeamertemplate{background canvas}{\includegraphics[width=\paperwidth,height=\paperheight]{fondo.png}}
\setbeamertemplate{footline}{}
\begin{frame}[plain]
    \vfill
    \begin{center}
        {\color{white}\Huge ¡Gracias!}

        \vspace{10pt}
        {\color{white}\large ¿Preguntas?}
    \end{center}
    \vfill
\end{frame}
}
